\section{Introduction}

% Instructions for writing the introduction section
%
% In the introduction section, there would be five subsection. Each section represents their own strength. The details of each five sections
% has been provided below
%
%
% Subsection 01 | Background and Motivation
%
% Explain the use of LLMs for Bangla and why language representation matters.
% Introduce subword tokenization and confidence calibration in accessible terms.
%
%
% Subsection 02 | Problem Statement
%
% State the problem precisely: whether item-level tokenization fragmentation is
% associated with model confidence/calibration behaviour on Bangla tasks.
% Distinguish accuracy (being correct) from calibration (confidence matching accuracy).
%
%
% Subsection 03 | Research Gap
%
% Support the gap with the reviewed literature. Do not claim that no prior work exists
% unless the literature review justifies that statement.
% State what prior tokenizer studies and Bangla LLM evaluations do not jointly test.
%
%
% Subsection 04 | Research Questions and Contributions
%
% RQ 01 : How does tokenization fragmentation differ between Bangla and English under the selected model tokenizers?
% RQ 02 : Is item-level tokenization fragmentation associated with confidence and calibration behaviour on Bangla evaluation items?
% RQ 03 : Does this association vary across models, subjects, or language conditions?
%
% List 2--4 concrete contributions only after the protocol and evidence are established.
%
%
% Subsection 05 | Paper Organization
%
% One short paragraph explaining the remaining sections.
%
%
% Please, make sure that, the introduction section contains the reference of the literature. Do not assume or include any information that is not
% cited or not taken from the literature.
